\clearpage
\fancyhead[L]{\small ASTRONOMIE.}
\newcommand{\JunoSymbol}{\text{\mdseries\AstroSymbols ⚵}}
\SourceLeaf{249}%

\Needspace{10\baselineskip}
\begin{center}
PALLAS.
\end{center}


Monatliche Correspondenz. Band VII. S.\ 369—375. S.\ 466—469. 1803 April u.\ Mai.


[Das April-Heft gibt den Inhalt des Briefes an \textsc{Bode} vom 3.\ März V.\ oben S.\ 233. Das Mai-Heft enthält noch folgende Stelle:]

— — — Diese Beobachtungen [des Dr.\ \textsc{Olbers}] stimmen noch vortrefflich mit den VI.\ Elementen der Bahn des Dr.\ \textsc{Gauss}; hier ist die Vergleichung:

\begin{center}
\begingroup\fontsize{10}{12}\selectfont\setlength{\tabcolsep}{3pt}\renewcommand{\arraystretch}{1.12}
\begin{tabular}{lrl}
\multicolumn{3}{l}{1803 März 4. $17^u\,11^m\,41^s$ m. Z. Br.} \\
Berechn. AR. $\PallasSymbol$ & $275^\circ\,52'\,35''9$ & Unterschied. R—B, $-2''1$ \\
Berechn. Decl. $\PallasSymbol$ & $8\ 58\ 18.1$ & Unterschied. R—B, $-4.9$ \\
\end{tabular}\endgroup
\end{center}

Da der Fehler der vorhergehenden V.\ Elemente so klein und ausserdem immer hinreichend genau bekannt ist, so wird die in unserm December-Hefte gegebene Ephemeride in ihrer ganzen Dauer zur Auffindung vollkommen hinreichen. Die gerade Aufsteigung darin ist jetzt $2\tfrac{1}{2}$ Min.\ zu gross und wird am 9.\ August $4'\,6''$ zu gross werden, in so fern man sich auf die fortwährende Übereinstimmung der neuen Element verlassen kann. — — — Dr.\ \textsc{Gauss} hat indessen [in der von ihm bisher angeordneten Form] die Ephemeride der \emph{Pallas} nach den neuen VI.\ Elementen [von Aug.\ 9] bis zum 23.\ Oct.\ 1803 fortgesetzt; länger wird man sie wol nicht verfolgen können, da sie schon den 18.\ Septemb.\ lichtschwächer zu werden anfängt, als sie bei ihrer Wiederentdeckung war.

\begin{flushright}
[\textsc{von Zach}.]
\end{flushright}

\Needspace{28\baselineskip}
\begin{center}
PALLAS.
\end{center}


Monatliche Correspondenz. Band VII. S.\ 556—560. 1803 Juni.


— — —

Dr.\ \textsc{Gauss} fährt unermüdet fort, die \textsc{Olbers}'schen Beobachtungen der \emph{Pallas} mit seinen VI.\ Elementen zu vergleichen; hier das Resultat der Observationen vom 16.\ März bis zum 13.\ April, woraus man mit Vergnügen ersehen wird, dass jene Elemente noch fortdauernd gut übereinstimmen, und dass die Unterschiede kaum grösser sind, als die bei den Beobachtungen selbst vorauszusetzenden Ungewissheiten, besonders in den Abweichungen, welche vom Dr.\ \textsc{Olbers} selbst als sehr unsicher angegeben werden:

\begin{center}
\begingroup\fontsize{8}{10}\selectfont\setlength{\tabcolsep}{3pt}\renewcommand{\arraystretch}{1.12}
\begin{tabular}{r|r|r|r|r|r}
1803 & \shortstack{Mittlere Zeit\\in\\Bremen} & \shortstack{Berechnete\\gerade Aufsteig.\\der Pallas} & \shortstack{Berechnete\\Abweichung der\\Pallas nördlich} & \shortstack{Unterschied\\in AR.\\R—B} & \shortstack{in Declin.\\R—B} \\ \hline
März 16 & $14^u\,10^m\,28^s$ & $278^\circ\,37'\,10''1$ & $10^\circ\,43'\,18''7$ & $-5''9$ & $+57''7$ \\
21 & $13\,40\,47$ & $279\,37\,19.0$ & $11\,30\,22.1$ & $+14.0$ & $+19.1$ \\
22 & $13\,14\,22$ & $279\,48\,30.2$ & $11\,39\,49.6$ & $+20.2$ & $+64.6$ \\
24 & $13\,2\,12$ & $280\,10\,29.9$ & $11\,59\,11.0$ & $+4.9$ & $+28.0$ \\
31 & $13\,24\,25$ & $281\,20\,15.1$ & $13\,8\,54.8$ & $-2.9$ & $+4.8$ \\
April 11 & $12\,10\,9$ & $282\,43\,33.5$ & $15\,0\,59.4$ & $-16.5$ & $+9.4$ \\
12 & $12\,12\,18$ & $282\,49\,32.1$ & $15\,11\,21.3$ & $-11.9$ & $+4.3$ \\
13 & $12\,13\,41$ & $282\,55\,13.0$ & $15\,21\,42.8$ & $-2.0$ & $+24.8$ \\
\end{tabular}\endgroup
\end{center}

\begin{flushright}
[\textsc{von Zach}.]
\end{flushright}
\Needspace{28\baselineskip}
\SourceLeaf{250}%
\begin{center}\textsc{Pallas.}\end{center}

\medskip
Monatliche Correspondenz. Band VIII. S. 90---93. 1803 Juli.

\medskip
— — —

\medskip
{\rightskip=0pt plus 1em
Dr. \textsc{Gauss} hat die Vergleichung seiner VI. Elemente mit den \textsc{Olbers}'schen Beobachtungen fortgesetzt, und folgende Übereinstimmung erhalten:\par}


\begin{center}
\begingroup\fontsize{9}{11}\selectfont\setlength{\tabcolsep}{3pt}\renewcommand{\arraystretch}{1.12}
\begin{tabular}{r|r|r|r|r|r}
1803 & \shortstack{Mittlere Zeit\\in\\Bremen} & \shortstack{Berechnete\\AR.\\der Pallas} & \shortstack{Berechnete\\Declination\\der Pallas} & \shortstack{Unterschied\\in AR.\\R—B} & \shortstack{in Declin.\\R—B} \\ \hline
April 15 & $12^u\,3^m\,34^s$ & $283^\circ\,5'\,39''9$ & $15^\circ\,42'\,20''3$ & $-\,1''1$ & $+\,9''3$ \\
20 & $12\,55\,28$ & $283\,26\,47.3$ & $16\,34\,18.4$ & $-\,8.7$ & $-\,32.6$ \\
24 & $11\,50\,56$ & $283\,38\,3.9$ & $17\,14\,30.4$ & $+\,20.9$ & $-\,16.6$ \\
25 & $12\,16\,41$ & $283\,40\,0.6$ & $17\,25\,0.7$ & $+\,15.6$ & $-\,30.3$ \\
Mai 11 & $11\,22\,17$ & $283\,29\,40.1$ & $19\,57\,28.6$ & $-\,12.9$ & $+\,26.6$ \\
12 & $11\,50\,10$ & $283\,26\,14.9$ & $20\,6\,23.9$ & $-\,15.1$ & $+\,19.9$ \\
20 & $11\,38\,40$ & $282\,48\,3.6$ & $21\,11\,25.3$ & $+\,4.6$ & $+\,4.3$ \\
31 & $10\,57\,47$ & $281\,24\,52.3$ & $22\,22\,35.8$ & $+\,2.3$ & $-\,10.2$ \\
Juni 1 & $10\,52\,5$ & $281\,15\,44.6$ & $22\,27\,49.1$ & $-\,8.4$ & $-\,8.9$ \\
1 & $11\,9\,33$ & $281\,15\,37.9$ & $22\,27\,52.9$ & $-\,1.3$ & $-\,13.1$ \\
\end{tabular}\endgroup
\end{center}

\begin{flushright}
[\textsc{von Zach}.]
\end{flushright}

\bigskip
\Needspace{10\baselineskip}
\centerline{\textsc{Ceres.}}

\medskip
Monatliche Correspondenz. B. VIII. S. 288---291 u. S. 369---371. 1803 Sept. u. Oct.

\medskip
— — —

\medskip
Dr. \textsc{Gauss} untersuchte indessen die, aus den Bremer und den drei im Julius-Stück S. 94 angegebenen Palermer Beobachtungen hervorgehende Correction seiner letzten (VIII) Elemente. Eine nur leichte Veränderung ist hinreichend gewesen, diese neuen Beobachtungen mit den alten zu vereinigen. Nur den Knoten musste er 3 Min. weiter rücken, als für die Beobachtungen in Palermo von 1801. Hoffentlich wird sich dies künftig bei vollständigerer Rechnung der Störungen rechtfertigen lassen. Hier sind indessen diese neuen (IX) Elemente.

\begin{center}
\begingroup\fontsize{9}{11}\selectfont\setlength{\tabcolsep}{3pt}
\begin{tabular}{lr}
Epoche für Seeberg 1803 . . . . . . . . . . . . & $233^\circ\ 36'\ 3''1$ \\
Tropische Bewegung in 365 Tagen . . . . . . . . . & $78\ 8\ 7.2$ \\
Tropische Bewegung in einem Tage . . . . . . . . & $770''650$ \\
Logarithmus der halben grossen Axe . . . . . . . . & $0.4421516$ \\
Sonnenferne 1803 . . . . . . . . . . . . . . . & $326^\circ\ 33'\ 18''$ \\
Knoten 1803 . . . . . . . . . . . . . . . . . & $80\ 58\ 22$ \\
Excentricität . . . . . . . . . . . . . . . . & $0.0788941$ \\
Neigung . . . . . . . . . . . . . . . . . . & $10^\circ\ 37'\ 54''$ \\
\end{tabular}\endgroup
\end{center}
\SourceLeaf{251}%

Diese Elemente stimmen mit den erwähnten drei \textsc{Piazzi}'schen Beobachtungen so:

\begin{center}
\begingroup\fontsize{9}{11}\selectfont\setlength{\tabcolsep}{3pt}\renewcommand{\arraystretch}{1.12}
\begin{tabular}{r|r|r|r|r|r}
1803 & \shortstack{Mittlere Zeit\\in Palermo} & \shortstack{Berechnete\\AR. der $\CeresSymbol$} & \shortstack{Abweich der $\CeresSymbol$} & \shortstack{Unterschied R—B\\in AR.} & \shortstack{in Decl.} \\ \hline
Mai 12 & $15^u\,53^m\,37^s7$ & $288^\circ\,19'\,17''9$ & $24^\circ\,36'\,22''0$ & $+\,2''9$ & $+\,9''0$ \\
--- 13 & $15\,49\,41.3$ & $288\,18\,57.3$ & $24\,39\,57.2$ & $-\,10.2$ & $-\,4.2$ \\
--- 14 & $15\,45\,41.9$ & $288\,18\,13.1$ & $24\,43\,36.7$ & $-\,1.9$ & $+\,5.6$ \\
\end{tabular}\endgroup
\end{center}

[Das October-Heft enthält die nach den IX. Elementen von \textsc{Gauss} berechnete Ephemeride von 1804 April 30. bis 1805 Januar 19 mit dreitägigen Zwischenräumen für Mitternacht in Seeberg, die Gerade Aufsteigung und Abweichung auf Minuten genau, die Lichtstärke in fünf Decimalstellen und auch die mittlere Zeit des Durchganges durch den Meridian.]

\begin{flushright}
[\textsc{von Zach}.]
\end{flushright}

\bigskip
\Needspace{10\baselineskip}
\centerline{\textsc{Ceres und Pallas.}}

\medskip
Monatliche Correspondenz. Band IX. S. 246---250. 1804. März.

\medskip
Wir haben schon im October-Hefte vorigen Jahrs S. 370 eine von Dr. \textsc{Gauss} nach seinen zum neuntenmal verbesserten Elementen der \emph{Ceres}-Bahn berechnete Ephemeride des geocentrischen Laufes dieses Planeten für das Jahr 1804 mitgetheilt und im December-Hefte S. 535 eine Karte dieses Laufes von dem Inspector \textsc{Harding} in Lilienthal versprochen. Diese Karte folgt nunmehr mit gegenwärtigem Hefte.

{\rightskip=0pt plus 1em
Dr. \textsc{Gauss} hat seitdem die Elemente der \emph{Pallas}-Bahn nach den spätesten Beobachtungen noch einmal corrigirt, und nach diesen zum VII. mal verbesserten Bestandtheilen dieser Bahn die nachstehende Ephemeride berechnet, welcher wir gleichfalls eine vom Inspector \textsc{Harding} entworfene Karte des geocentrischen Laufes dieses Planeten nachfolgen lassen werden. [Die Ephemeride geht von 1804 April 30 bis 1805 Januar 19 und gibt in der bisher angewandten Form Gerade Aufsteigung, Abweichung und Entfernung von der Erde und Lichtstärke.]\par}


An einigen Stellen ist der Unterschied von der vom Prof. \textsc{Bode} in dem Berl. Astr. J. B. 1806 S. 91 gelieferten Ephemeride beträchtlich. Die Beobachtungen der \emph{Pallas} werden in diesem Jahre nun schon wieder etwas leichter sein, als im vorigen. Die grösste Lichtstärke vom 20. Juni vor. J. erreicht sie in diesem Jahre am 23. Juli und wird am 7. Sept. noch mehr als $\frac{1}{4}$ heller, wozu noch die völlige Dunkelheit der Nächte und der Umstand kommt, dass der \emph{Pegasus} und \emph{Wassermann} nicht so sternreich sind, als die Gegend, worin der Planet sich voriges Jahr aufhielt. Die Lichtstärke, bei der er voriges Jahr zuerst aufgefunden wurde, erreichte er diesesmal schon am 5. Juni und am 18. Mai die vom 10. October 1803, bei der Dr. \textsc{Olbers} ihn zum letztenmal beobachtete.

Die neuesten und VII. Elemente, nach welchen Dr. \textsc{Gauss} obige Ephemeride der \emph{Pallas}-Bahn berechnet hat, sind folgende:

\begin{center}
\begingroup\fontsize{9}{11}\selectfont\setlength{\tabcolsep}{3pt}
\begin{tabular}{lr}
Epoche (Seeberg) 1803 . . . . . . & $221^\circ\ 29'\ 32''0$ \\
tägliche tropische Bewegung . . . & $770''0446$ \\
jährliche tropische Bewegung . . . & $78^\circ\ 4'\ 26''3$ \\
Sonnenferne 1803 . . . . . . . & $301\ 17\ 34.4$ \\
Excentricität . . . . . . . . . & $0.2457396$ \\
Logarithmus der halben grossen Axe  & $0.4423790$ \\
Aufsteig. Knoten 1803 . . . . . . & $172^\circ\ 28'\ 13''7$ \\
Neigung der Bahn . . . . . . . & $34\ 38\ 1.1$ \\
\end{tabular}\endgroup
\end{center}
\SourceLeaf{252}%

So wie wir im letzten December-Hefte S. 536 ein kleines Verzeichniss derjenigen Sterne, welche sich auf dem Wege der \emph{Ceres} oder in ihrer Nähe befinden werden, aus \textsc{Piazzi}'s grossem Stern-Verzeichniss mitgetheilt haben, eben so geben wir auch hier einen kleinen Auszug derjenigen Sterne aus diesem Verzeichniss, welche mit der \emph{Pallas} sehr nahe zusammenkommen, und mit ihr verglichen werden können.

\begin{flushright}
[\textsc{von Zach}.]
\end{flushright}

\bigskip
\Needspace{10\baselineskip}
\centerline{\textsc{Neuer Comet.}}

\medskip
Monatliche Correspondenz. Band IX. S. 432---435. 1804 Mai.

\medskip
Der von Dr. \textsc{Olbers} den 12. März entdeckte neue Comet wurde auch zu Marseille, auf der Sternwarte der Marine, den 7. März von \textsc{Pons} entdeckt. Den 10. März entdeckte ihn auch der Astronom der Pariser National-Sternwarte \textsc{Bouvard}, und fand um $15^u\ 57^m$ m. Z. dessen gerade Aufsteigung $220^\circ\ 4'$, dessen südliche Abweichung $1^\circ\ 41'$; den 11. März um $15^u\ 39^m$ m. Z., gerade Aufsteigung $220^\circ\ 10'$, nördl. Abweichung $2^\circ\ 44'$; den 29. März um $9^u$ m. Z., gerade Aufsteig. $218^\circ\ 40'$, nördl. Abweichung $46^\circ\ 37'$.

Auch \textsc{Messier} verfolgte diesen Cometen in Paris, und beobachtete den 17. März um $11^u\ 34^m\ 18^s$ w. Z. gerade Aufsteigung $220^\circ\ 29'\ 45''$, nördl. Abweichung $26^\circ\ 18'\ 24''$.

\Needspace{18\baselineskip}
Dr. \textsc{Olbers} beobachtete diesen Cometen zehn Tage bis zum 1. April. Aus diesen Beobachtungen hat Dr. \textsc{Gauss} folgende parabolische Elemente seiner Bahn herausgebracht:

\begin{center}
\begingroup\fontsize{9}{11}\selectfont\setlength{\tabcolsep}{3pt}
\begin{tabular}{lr}
Zeit der Sonnennähe 1804 Febr. 13. . . . . . & $14^u\ 49^m\ 51^s$ m. Seeb. Z. \\
Logarithmus des Abstandes . . . . . . . . & $0.0298575$ \\
\end{tabular}\endgroup
\end{center}
{\fontsize{8}{10}\selectfont
\[\left.\begin{array}{ll}\text{Länge der Sonnennähe . . . . . . . . . .}&148^\circ\ 44'\ 51''\\\text{Länge des aufsteigenden Knotens . . . . . .}&176\ 47\ 58\end{array}\right\}\begin{array}{l}\text{beide vom mittl.}\\\text{Aequinoctium gezählt}\end{array}\]}
\begin{center}
\begingroup\fontsize{9}{11}\selectfont\setlength{\tabcolsep}{3pt}
\begin{tabular}{lr}
Neigung der Bahn . . . . . . . . . . . & $56\ 28\ 40$ \\
Bewegung . . . . . . . . . . . . . . & rechtläufig. \\
\end{tabular}\endgroup
\end{center}

Nach diesen Elementen steht die Vergleichung mit den sämmtlichen \textsc{Olbers}'schen Beobachtungen folgendermassen:

\begin{center}
\begingroup\fontsize{9}{11}\selectfont\setlength{\tabcolsep}{3pt}\renewcommand{\arraystretch}{1.12}
\begin{tabular}{r|r|r|r}
1804 & \shortstack{Mittl. Zeit\\Bremen} & \shortstack{Unterschied der Rechnung\\und Beobachtung\\in AR.} & \shortstack{in Decl.} \\ \hline
März 12 & $12^u\,56^m\,13^s$ & $+9''$ & $0''$ \\
13 & $11\,40\,43$ & $-19$ & $-7$ \\
14 & $12\,22\,26$ & $+22$ & $+92$ \\
15 & $8\,54\,41$ & $-31$ & $-79$ \\
20 & $9\,22\,52$ & $-247 ::$ & $+121 ::$ \\
\end{tabular}\endgroup
\end{center}
\begin{center}
\begingroup\fontsize{9}{11}\selectfont\setlength{\tabcolsep}{3pt}\renewcommand{\arraystretch}{1.12}
\begin{tabular}{r|r|r|r}
1802 & \shortstack{Mittl. Zeit\\Bremen} & \shortstack{Unterschied der Rechnung\\und Beobachtung\\in AR.} & \shortstack{in Decl.} \\ \hline
März 22 & $8^u\,59^m\,13^s$ & $-143''$ & $+85''$ \\
27 & $8\,59\,43$ & $-55$ & $-152$ \\
28 & $8\,28\,2$ & $-27$ & $+32$ \\
29 & $8\,45\,41$ & $-26$ & $+33$ \\
April 1 & $9\,1\,52$ & $+28$ & $-1$ \\
\end{tabular}\endgroup
\end{center}

\noindent wo die Fehler besonders in der Declination so irregulair laufen, dass man sie füglich grossentheils den Beobachtungen zuschreiben kann. Die Beobachtung vom 20. März wird von Dr. \textsc{Olbers} selbst als ganz zweifelhaft angegeben.

Bei dieser Vergleichung hat Dr. \textsc{Gauss} sowol auf. Nutation und Aberration als auf die Parallaxe Rücksicht genommen; letztere ist besonders bei den ersten Beobachtungen, wo der Abstand des Come%
\SourceLeaf{253}%
ten von der Erde nur ein Viertel der Entfernung der Sonne war, nicht ganz unerheblich; er hat sie daher um so lieber mit in Betrachtung gezogen, da die Arbeit dadurch fast gar nicht vergrössert wird.
Er hat sich nemlich bei diesen Rechnungen des Verfahrens, welches in dem im gegenwärtigen Hefte S. 385 [S. 94 d. B.] abgedruckten vortrefflichen Aufsatze erklärt worden, bedient, und für die Constanten $A$, $B$, $C$, $a$, $b$, $c$ folgende Werthe zur unmittelbaren Berechnung der geraden Aufsteigung und Declination gefunden:
\[
A = 268^\circ\ 14'\ 21'', \quad B = 176^\circ\ 29'\ 48'', \quad C = 2^\circ\ 19'\ 58''
\]
\[
a = 87\ \ 19\ \ 38, \quad b = 56\ \ 56\ \ 7, \quad c = 33\ \ 12\ \ 3
\]

\Needspace{22\baselineskip}
\indent Die \emph{Constanten} beziehen sich schon auf den wahren Aequator, und schliessen also die während einer so kurzen Zeit als unveränderlich anzusehende Nutation mit ein; ist also $v$ die wahre Anomalie des Cometen, so werden die Coordinaten $x$, $y$, $z$ sogleich durch folgende höchst einfache Formeln gefunden:

\[
\begin{aligned}
x &= \frac{\alpha \sin(v + 240^\circ\ 11'\ 14'')}{\cos\frac{1}{2}v^2}, \\[5pt]
y &= \frac{\ObservedBeta \sin(v + 148^\circ\ 26'\ 41'')}{\cos\frac{1}{2}v^2}, \\[5pt]
z &= \frac{\gamma \sin(v + 334^\circ\ 16'\ 51'')}{\cos\frac{1}{2}v^2}
\end{aligned}
\]

\noindent und wo

\[
\log\alpha = 0.029384, \quad \log \ObservedBeta = 9.953130, \quad \log \gamma = 9.768302
\]

\indent Die Grössen $\alpha$, $\ObservedBeta$, $\gamma$ sind hier nemlich die Producte aus dem Abstande im Perihelium mit dem Sinus von $a$, $b$, $c$; und die Winkel $240^\circ\ 11'\ 14''$ u.\,s.\,w. gleich den Winkeln $A$, $B$, $C$ weniger der Entfernung des Knotens von der Sonnennähe $=\,28^\circ\ 3'\ 7''$. — — —

\begin{flushright}
[\textsc{von Zach}.]
\end{flushright}

\bigskip
\Needspace{10\baselineskip}
\centerline{\textsc{Gauss an Bode.}}

\medskip
Astronomisches Jahrbuch für 1808. S. 187---190. Berlin 1808.

\medskip
\begin{flushright}
Braunschweig 1804 Sept. 30.
\end{flushright}

Ew. --- nehme ich mir die Freiheit hier meine sämmtlichen bisherigen Beobachtungen der Juno vorzulegen. Ich habe sie, zwar nur am Kreismikrometer, aber mit möglichster Sorgfalt angestellt. Alle sind die Resultate aus sehr zahlreichen Vergleichungen theils mit Sternen aus \textsc{Piazzi}'s Catalog, theils mit solchen aus der Hist. Céleste, die ich nach diesem Catalog sorgfältig reducirt habe. Die 3 ersten ausgenommen, welche mit einem schlechtern Werkzeuge gemacht sind, wurden sie sämmtlich mit einem sehr guten \textsc{Short}'schen Spiegelteleskope angestellt.

\begin{center}
\begingroup\fontsize{10}{12}\selectfont\setlength{\tabcolsep}{3pt}\renewcommand{\arraystretch}{1.12}
\begin{tabular}{r|rrr|rrr|rrr}
1804 & \multicolumn{3}{c|}{Mittl. Zeit} & \multicolumn{3}{c|}{\shortstack{Scheinbare\\Gerade Aufst.}} & \multicolumn{3}{c}{Südl. Decl.} \\ \hline
Sept. 12 & $10^u$ & $35^m$ & $2^s$ & $0^\circ$ & $54'$ & $31''$ & $1^\circ$ & $38'$ & $15''$ \\
13 & $9$ & $41$ & $32$ &  & $45$ & $24$ &  & $50$ & $59$ \\
14 & $11$ & $38$ & $49$ &  & $35$ & $37$ & $2$ & $5$ & $1$ \\
15 & $10$ & $16$ & $17$ &  & $26$ & $53$ &  & $17$ & $35$ \\
16 & $10$ & $37$ & $4$ &  & $17$ & $17$ &  & $31$ & $20$ \\
17 & $11$ & $28$ & $59$ &  & $7$ & $23$ &  & $44$ & $29$ \\
\end{tabular}\endgroup
\end{center}
\begin{center}
\begingroup\fontsize{10}{12}\selectfont\setlength{\tabcolsep}{3pt}\renewcommand{\arraystretch}{1.12}
\begin{tabular}{r|rrr|rrr|rrr}
1804 & \multicolumn{3}{c|}{Mittl. Zeit} & \multicolumn{3}{c|}{\shortstack{Scheinbare\\Gerade Aufst.}} & \multicolumn{3}{c}{Südl. Decl.} \\ \hline
Sept. 18 & $11^u$ & $22^m$ & $16^s$ & $359^\circ$ & $57'$ & $26''$ &  & $58'$ & $3''$ \\
21 & $10$ & $24$ & $52$ &  & $28$ & $6$ & $3$ & $38$ & $0$ \\
24 & $10$ & $1$ & $2$ & $358$ & $57$ & $53$ & $4$ & $18$ & $6$ \\
25 & $8$ & $44$ & $25$ &  & $48$ & $12$ &  & $30$ & $44$ \\
27 & $10$ & $20$ & $29$ &  & $27$ & $20$ &  & $57$ & $47$ \\
28 & $8$ & $29$ & $4$ &  & $18$ & $20$ & $5$ & $10$ & $22$ \\
\end{tabular}\endgroup
\end{center}
\SourceLeaf{254}%

Theils nach diesen Beobachtungen, theils nach denen unsers gemeinschaftlichen Freundes
Hr. Dr. \textsc{Olbers} und einigen des Frh. \textsc{von Zach} habe ich bereits eine erste Annäherung zu
den Elementen versucht, die ich die Ehre habe Ihnen hier mitzutheilen. Ich halte die
Planeten-Natur der Juno nunmehr bereits \emph{aus den Beobachtungen} für erwiesen und
schmeichle mir, dass auch die Elemente schon eine Idee von den genäherten Dimensionen
der Bahn zu geben zureichend sind,

\begin{center}
I\textsuperscript{te} Elemente der Juno-Bahn.
\end{center}

\begin{center}
\begingroup\fontsize{8}{10}\selectfont\setlength{\tabcolsep}{3pt}
\begin{tabular}{lr}
Epoche 1804. Sept. 5 $0^u$ M. Z. im Seeberger Meridian . & $20^\circ\ 38'\ 56''$ \\
tägliche mittlere Bewegung . . . . . . . . . . & $779''80$ \\
Sonnenferne . . . . . . . . . . . . . . . . & $239^\circ\ 14'\ 2''$ \\
Logar. der halben Axe . . . . . . . . . . . . . & $0.438682$ \\
Excentricität . . . . . . . . . . . . . . . & $0.287359$ \\
Aufsteigender Knoten . . . . . . . . . . . . . . & $171^\circ\ 15'\ 35''$ \\
Neigung . . . . . . . . . . . . . . . . . . & $13\ 34\ 59$ \\
\end{tabular}\endgroup
\end{center}

Natürlich werden diese Elemente, da die bisherige Dauer der Beob. noch so kurz ist,
noch ansehnlicher Verbesserungen bedürfen, und namentlich ist die Ungewissheit bei der
mittl. Bewegung noch viel grösser als der kleine Unterschied von der mittl. Bewegung der
$\CeresSymbol$ und $\PallasSymbol$, daher es gar wol sein kann, dass alle 3 der Gleichheit noch viel näher kommen.

Um den entferntern Astronomen, denen etwa die Auffindung des Planeten noch nicht gelungen
ist, dieselbe nach meinem Vermögen zu erleichtern, habe ich nach obigen Elementen eine
kleine Ephemeride berechnet, die hoffentlich den Lauf desselben auf mehrere Wochen
hinreichend genau zur Auffindung darstellen wird.

\Needspace{26\baselineskip}
\begin{flushright}
Braunschweig 1804 November 12.
\end{flushright}

Indem ich Ew.--- für die mir von Ihnen mitgetheilten Beobachtungen meinen
verbindlichsten Dank abstatte, fahre ich zugleich fort, Ihnen von den letzten
Resultaten meiner Untersuchungen über die Bahn der Juno Nachricht zu geben, obzwar
dieselben bereits vor 14 Tagen berechnet sind, und also bald einer neuen Verbesserung
fähig sein werden. Hier die

\begin{center}
III\textsuperscript{ten} Elemente der Juno.
\end{center}

\begin{center}
\begingroup\fontsize{9}{11}\selectfont\setlength{\tabcolsep}{3pt}
\begin{tabular}{lr}
Epoche 1805 im Mer. von Seeberg . . . . . . . . & $43^\circ\ 21'\ 2''$ \\
Sonnenferne . . . . . . . . . . . . . . . . . & $233\ 56\ 6$ \\
Neigung . . . . . . . . . . . . . . . . . . & $12\ 52\ 48$ \\
Logarithmus der halben Axe . . . . . . . . . . & $0.426699$ \\
tägliche Bewegung . . . . . . . . . . . . . . & $812''754$ \\
Knoten . . . . . . . . . . . . . . . . . . & $171^\circ\ 0'\ 0''$ \\
Excentricität . . . . . . . . . . . . . . . . & $0.263182$ \\
\end{tabular}\endgroup
\end{center}

Dass die mittlere Bewegung der Juno der der $\CeresSymbol$ und $\PallasSymbol$
nicht gleich, sondern grösser sei, sehe ich bereits als ausgemacht an.

Bei Berechnung der III\textsuperscript{ten} Elemente hatte ich die Beob. des Hrn. \textsc{von Zach} bis 21. Oct. vor mir, mit denen sie bis auf ein paar Secunden übereinstimmen.
Dr. \textsc{Maskelynes} Beobachtungen habe ich bisher bis zum 17\textsuperscript{ten}. Nach den
spätern Beobachtungen aber ist es entschieden, dass diese Elemente die AR. jetzt zu
gross geben, wegen der Decl. aber kann ich noch nichts entscheiden. Hier die
Vergleichung der neuesten Beobachtungen.
\SourceLeaf{255}%


\begin{center}\begin{tabular}{r|cc|c|cc}
 & \multicolumn{2}{c|}{Hr. Prof. \textsc{Bode}} & & \multicolumn{2}{c}{Hr. Dr. \textsc{Olbers}} \\
 & AR. & Decl. & & AR. & Decl. \\
\hline
Oct. 20 & $+33''$ & $+7''$ & Oct. 27 & $+21''$ & $+15''$ \\
Nov.\ 2 & $+48$ & $+5$ & 30 & $+34$ & $+12$ \\
5 & $+46$ & $-15$ & Nov.\ 2 & $+33$ & $-1$ \\
 & & & 5 & $+39$ & $+25$ \\
 & & & 6 & $+37$ & $+32$ \\
\end{tabular}\end{center}

Den 20. Dec. kommen $\CeresSymbol$ und $\JunoSymbol$ so nahe zusammen, dass man
sie zugleich im Fernrohr wird beobachten können.


